% ======================= 第 八 章 =======================
\chapter{撤离与落地}
\chaptersubtitle{ГЛАВА ВОСЬМАЯ}{ЭВАКУАЦИЯ И ПРИЗЕМЛЕНИЕ}
\epigraphbox{“天作孽，犹可违；自作孽，不可活。”}{《尚书·太甲》}

\kaishou{本章要说的是：这座城不是被打败的，是修不下去了。前二二一年它把最后一批人放下去，此后月球上再没有月之民。}

\noindent
第七章末尾那个数字要在这里兑现。月兔之乱以后，城又运转了一千七百六十年。这一千七百年里，系统没坏，人也没死，只是每年都要多花一点力气去压制山里的东西。

压到某一年，压不动了。

\section{远因、近因与直接原因}

撤离的原因是三层，档案里只记最后一层。

远因是污秽的积累。月兔死尽之后，入库的是原状污秽，不稳定，会扩散。系统仍在抽取，抽出来的东西无处加工，只能压在山里。一千七百六十年间，山里的压力一直在涨，没有回落。

近因是能量。月之都的供能、控温、生命保障，都依附于地球这颗行星的热机。地球内部的对流效率在缓慢下降，这一点在前四二一年与前二二一年两次测站校准中测得，二百年间下降百分之三点一。

直接原因是前二二一年的生态崩溃。抽取污秽的那一段停转，城内污秽失控。到这一步，已经没有别的办法。

\begin{table}[htbp]
\centering
\caption{六批撤离}
\label{tab:six}
\small
\begin{tabular}{@{}clrl@{}}
\toprule
批 & 年份 & 人数 & 性质\\
\midrule
一 & 前 421 & 6,000 & 计划内，先行探路\\
二 & 前 380 & 8,000 & 计划内\\
三 & 前 340 & 10,000 & 计划内\\
四 & 前 300 & 13,000 & 计划内\\
五 & 前 260 & 18,000 & 计划内，末次正常轮换\\
六 & 前 221 & 65,000 & 计划外，应急\\
\midrule
 & 合计 & 120,000 & \\
\bottomrule
\end{tabular}
\tabsource{据第 I 卷简二十九至三十一与本书第六章所引档案。前五批的年份在第 I 卷中只给到“某年”，此处所补为其推算值。}
\end{table}

表 \ref{tab:six} 的最后一行与前五行不是一套账。实载逐批记数，合计十二万；而在册的记法是另一样：整次撤离只记一个总数，十一万九千。

两个数差一千。

这一千人的去向，档案不说。第 I 卷推测是登记与登船之间的脱节，本书同意这个推测，不再另立新说。

\section{玉山再入与落地四事}

前二二一年，最后一批六万五千人乘玉山返回地面。

玉山是一艘再入载具，由一个无人主体与十二个乘员舱组成。主体全长四千七百一十三米，落地前解体，主体硬着陆于黄河流域，乘员舱分散落入豫西与关中之间。落点据测算为北纬三十四点六度、东经一百一十点四度。

落地之后发生的事，按顺序有四件。

第一件是沾秽。这是两万年来第一次。他们的身体在这两万年里没有接触过任何变化中的东西，落地即接触，接触即开始。

第二件是变化恢复。沾秽之后，无秽状态解除，他们开始衰老。这个过程不快，几十年里一点一点发作，与常人无异。

第三件是生育恢复。第一代之后出现了子女，子女同为凡人。

第四件是遗忘恢复。随身的那点残余效应重新发作，此前被冻住的判断开始松动。这一条最晚显现，也最容易被忽略。

四代之内，人口从十二万降到不足五千。存活率依次为百分之六十一点四、三十八点七、十二点九、四点二。

\begin{figure}[htbp]
\centering
\begin{tikzpicture}[x=1mm,y=1mm]
  \useasboundingbox (-4,-16) rectangle (140,20);
  \draw[ink,line width=0.4pt] (0,0) -- (132,0);
  \draw[ink,line width=0.4pt] (0,0) -- (0,16);
  % 四代
  \foreach \x/\h/\lab/\pct in {18/14.7/一代/61.4, 46/9.3/二代/38.7, 74/3.1/三代/12.9, 102/1.0/四代/4.2}{
    \fill[rulegrey] (\x-8,0) rectangle (\x+8,\h);
    \draw[ink,line width=0.4pt] (\x-8,0) rectangle (\x+8,\h);
    \node[anchor=north,font=\scriptsize] at (\x,-1.4) {\lab};
    \node[anchor=south,font=\scriptsize] at (\x,\h+0.4) {\pct\%};
  }
  \node[anchor=east,font=\scriptsize,align=right] at (-2,8) {人口\\（千人）};
  \foreach \y in {4,8,12}{
    \draw[rulegrey,line width=0.2pt] (0,\y) -- (132,\y);}
\end{tikzpicture}
\caption{落地后四代的人口}
\figsource{据第 I 卷简二十六所记存活率重绘。纵轴以初代十二万为百分之百折成，非实际人数。}
\label{fig:four}
\end{figure}

图 \ref{fig:four} 里第四代那根柱子只有一点，不到五千人。此后档案不再有数字。

四代之后，「华夏神族」作为一个群体消失了。消失的是认同，与血统无关：这一支的血统延续了千余年，唐宋以后才无法分辨。

\begin{archive}[机密／СЕКРЕТНО]{档案 № ФВ-1973-08}
玉山再入记录，共二十三叶。

第 4 叶为乘员舱分布图，十二舱标号 A 至 L，各注落点范围。F 舱与 J 舱的落点范围被划掉，改写为「未寻获」。

第 11 叶记主体解体高度：海拔十一公里。解体原因记为「结构疲劳」，未附分析。

第 17 叶有一行铅笔字，字迹潦草，内容为：「他们下来的时候还有说有笑。」此句无署名，前后无关联文字。
\end{archive}

档案里那句没有署名的铅笔字，是本书所见关于落地最直接的一处人性记录。它不解释任何事，但它是唯一一处写到「笑」的地方。

\section{撤离的性质}

这一节要处理一个说法上的分歧。

前二二一年的撤离，在外的叙述里常被写成一次决断：这座城不值得再修，于是放弃。本书不取这个说法。

理由是材料。简二十七记着「再断，则会计」四个字，前一句是「三断三复」。这句话记的是一次下令：气路断了三次又修好三次，主事的人说，再断就清点人口。他下的令是清点，不是停修。

停修这件事，档案里没有。有的是生态崩溃：抽取污秽的一段停转，城内污秽失控。到这一步，修与不修已经不是一个可选项。

所以本书的说法是：撤离是被迫的。远因、近因、直接原因三层都在，没有一层是「他们决定不要这座城了」。

\section{落地十一年后的判断}

落地十一年后，一部分人判断中原不宜久居。这个判断有三条理由：秦始皇正在求仙，而他们就是仙；撤离前关于四代存活率的推算正在成为事实；沾秽之后他们从此会老。

这个判断的结果是东渡，见下一章。

留在中原的那一部分，接管并改造了既有聚落，没有另择新址。仰韶、龙山、二里头三处遗址的后期改造痕迹，与此可以对照。人口的主体始终是本土居民，十二万不足以立国。

秦始皇本人是土著。这一点在《战争的雷霆》第一章已经说清，本书不重复。

\section{本章结论}

本章说了四件事。

撤离有三层原因，档案只记最后一层。远因是污秽积累，近因是能量衰减，直接原因是生态崩溃。

六批合计十二万，在册十一万九千，差一千，去向不明。

落地四件事：沾秽、变化恢复、生育恢复、遗忘恢复。四代之内，群体消失。

撤离是被迫的。

本章的自我拆台：关于「落地之后」，本书没有任何一手材料。四代存活率是数字，不是叙述；没有一份月之民自己写的落地记录，也没有一份秦人写他们的记录。这一章的后半段全部是从数字与地层推的，推到「他们会老」为止，再往前一步就没有依据了。

笔者停在这里。再往下写，写的就是小说。
